\section{CVF-AE 方法}
\label{sec:method}

本文提出的 CVF-AE 方法在自适应进化搜索框架中引入约束可行性场（Constraint Viability Field, CVF）。其核心思想不是改变目标函数，也不是在评价后简单修复路径，而是在候选解生成阶段利用障碍物、禁飞区、高度、转弯、边界和风险信息形成一个有界的方向增量，使控制点在被评价之前就朝更可行、更稳健的区域移动。这样可以减少搜索在不可行区域内的无效迭代，同时保留进化算法对多目标质量折中的全局探索能力。

\subsection{总体框架}

CVF-AE 的一次迭代由基础进化搜索、约束状态识别、CVF 候选生成、稀疏局部修复和可行性优先接收组成。给定种群 $\{\mathbf{X}_i\}_{i=1}^{N}$，算法首先根据式 \eqref{eq:objective} 和式 \eqref{eq:total-violation} 评价每个个体，记录当前最优可行解、种群可行率、平均违反量和近期改进情况。随后，基础自适应进化算子生成候选方向 $\mathbf{D}_{\mathrm{AE}}$。当个体满足 CVF 触发条件时，算法额外构造约束可行性方向 $\mathbf{D}_{\mathrm{CVF}}$ 和质量引导方向 $\mathbf{D}_{\mathrm{Q}}$，并生成融合候选。

候选生成可写为
\begin{equation}
    \mathbf{X}_{\mathrm{AE}} =
    \Pi_{\Omega}
    \left(
    (1-\mu_t)(\mathbf{X}_i+\mathbf{D}_{\mathrm{AE}})
    +\mu_t\mathbf{X}_{\mathrm{elite}}
    \right),
    \label{eq:ae-candidate}
\end{equation}
\begin{equation}
    \begin{aligned}
    \mathbf{D}_s
    &=
    \alpha_s \mathbf{D}_{\mathrm{AE}}
    + \beta_s \mathbf{D}_{\mathrm{CVF}}\\
    &\quad
    + \gamma_s \mathbf{D}_{\mathrm{Q}},\\
    \mathbf{X}_{\mathrm{CVF}}
    &=
    \Pi_{\Omega}
    \Bigl(
    (1-\mu_t)(\mathbf{X}_i+\mathbf{D}_s)\\
    &\quad
    +\mu_t\mathbf{X}_{\mathrm{elite}}
    \Bigr),
    \end{aligned}
    \label{eq:cvf-candidate}
\end{equation}
其中 $\mathbf{D}_s$ 为状态融合方向，$\Pi_{\Omega}(\cdot)$ 表示边界投影，$\mathbf{X}_{\mathrm{elite}}$ 为当前精英池中采样得到的精英个体，$\mu_t$ 为随迭代推进的小权重锚定系数，$\alpha_s,\beta_s,\gamma_s$ 为由约束状态 $s$ 决定的融合权重。CVF 候选不是每代对所有个体都生成，而是通过稀疏触发控制额外评价次数。需要强调的是，CVF 分支是基础 AE 候选之外的额外候选，而不是全种群强制更新；CVF 候选必须通过局部接收规则，不能只依靠降低违反量进入种群；局部修复属于候选后处理，不等同于 CVF 场本身。

基础 AE 方向 $\mathbf D_{\mathrm{AE}}$ 按当前状态选择四类搜索模式。令 $\tau=t/T$，$\mathbf d_{\mathrm{best}}=\mathbf X^*-\mathbf X_i$，$\mathbf d_{\mathrm{elite}}=\bar{\mathbf X}_{\mathrm{elite}}-\mathbf X_i$，$\mathbf d_{\mathrm{sample}}=\mathbf X_{\mathrm{elite}}-\mathbf X_i$，$\mathbf d_{\mathrm{ref}}=\mathbf X_{\mathrm{ref}}-\mathbf X_i$；$\mathbf d_{ab}$ 和 $\mathbf d_{cd}$ 为从种群中随机抽取个体形成的两个差分向量，$\boldsymbol\xi$ 为按变量范围缩放的均匀随机扰动，$\boldsymbol\zeta$ 为高斯扰动。本文实现中的基础方向采用统一模板
\begin{equation}
    \begin{aligned}
    \mathbf D_{\mathrm{AE}}
    ={}&c_b\mathbf r_b\odot\mathbf d_{\mathrm{best}}
    +c_e\mathbf r_e\odot\mathbf d_{\mathrm{elite}}\\
    &+c_s\mathbf r_s\odot\mathbf d_{\mathrm{sample}}
    +c_r\mathbf r_r\odot\mathbf d_{\mathrm{ref}}\\
    &+c_1\mathbf d_{ab}+c_2\mathbf d_{cd}
    +c_u\boldsymbol\xi+c_g\boldsymbol\zeta .
    \end{aligned}
\label{eq:ae-state-direction}
\end{equation}
其中各 $\mathbf r$ 为 $[0,1]^{3K}$ 上的独立随机向量，$\odot$ 表示逐维乘积。四个状态分别对应可行性形成、可行性保持、质量细化和停滞恢复，其非零系数见表 \ref{tab:base-ae-coefficients}。

\begin{table}[tbp]
\centering
\caption{基础 AE 方向在不同状态下的非零系数}
\label{tab:base-ae-coefficients}
\scriptsize
\resizebox{\columnwidth}{!}{%
\begin{tabular}{@{}c c c c c c c c c@{}}
\toprule
$s$ & $c_b$ & $c_e$ & $c_s$ & $c_r$ & $c_1$ & $c_2$ & $c_u$ & $c_g$\\
\midrule
1 & 0 & $0.20+0.20(1-\tau)$ & 0 & $0.25+0.25(1-\tau)$ & 0.18 & 0 & $0.08+0.10(1-\tau)$ & 0\\
2 & $0.25+0.10\tau$ & 0.28 & 0 & 0.08 & 0.16 & 0 & 0.08 & 0\\
3 & $0.38+0.16\tau$ & 0 & $0.30+0.12\tau$ & 0 & 0.08 & 0 & 0.04 & 0\\
4 & 0.22 & 0 & 0 & 0 & 0.24 & 0.18 & 0 & 0.10\\
\bottomrule
\end{tabular}%
}
\end{table}

该基础方向与 CVF 方向分开定义，使消融实验中的 AE 基线和去除 CVF 版本具有明确复现口径。

\subsection{约束状态识别}

CVF-AE 根据种群当前搜索状态动态调整约束引导强度。令 $\rho_f$ 表示种群可行率，$\bar{V}$ 表示平均违反量，$D_{\mathrm{pop}}$ 表示种群多样性，$\Delta J_{\mathrm{recent}}$ 表示近期目标改进幅度：
\begin{equation}
    \rho_f=\frac{1}{N}\sum_{i=1}^{N}\mathbb I(\mathbf X_i\ \mathrm{feasible}),
    \qquad
    \bar V=\frac{1}{N}\sum_{i=1}^{N}V(\mathbf X_i),
    \label{eq:state-feasibility}
\end{equation}
\begin{equation}
    D_{\mathrm{pop}}=
    \frac{1}{3K}\sum_{j=1}^{3K}
    \mathrm{std}\left(\frac{X_{1j},\ldots,X_{Nj}}{u_j-l_j}\right),
    \label{eq:state-diversity}
\end{equation}
\begin{equation}
    \Delta J_{\mathrm{recent}}=
    \frac{\max(0,J^*_{t-w}-J^*_{t})}{\max(1,|J^*_{t-w}|)} ,
    \qquad w=15 .
    \label{eq:recent-improvement}
\end{equation}
当 $\Delta J_{\mathrm{recent}}<10^{-4}$ 时认为搜索停滞。观测状态按以下互斥顺序判定：
\begin{equation}
s_{\mathrm{obs}}=
\begin{cases}
4, & \begin{aligned}[t]
        &\Delta J_{\mathrm{recent}}<10^{-4}\\[-1mm]
        &{}\land (D_{\mathrm{pop}}<0.08\lor t>2w),
    \end{aligned}\\
1, & \rho_f<0.20\lor \bar V>10,\\
2, & \rho_f<0.65,\\
3, & \mathrm{otherwise}.
\end{cases}
\label{eq:observed-state}
\end{equation}
式 \eqref{eq:observed-state} 给出的是观测状态模板。正式方法进一步采用保守状态调度：默认处于可行性保持状态；仅当当前最优解尚不可行且 $\rho_f\leq0.10$ 或 $\bar V\geq18$ 时，请求切换到可行性形成状态；仅当搜索停滞、$\rho_f<0.55$ 且 $\bar V\geq6$ 时，请求切换到停滞恢复状态。非保持状态需要连续确认 $2$ 代才生效。质量细化状态保留为观测状态和权重模板；在本文保守调度下，若不满足形成或恢复切换条件，正式运行回到保持状态，因此本文不将质量细化写成单独决定性能的核心状态。该滞回规则减少了频繁状态跳变，避免 CVF 在已接近可行的阶段过度干预。

\subsection{约束可行性场构建}

CVF 在内部控制点空间中构造。实际构造时使用一个较低密度的路径采样集 $\widetilde{\mathcal P}=\{\widetilde{\mathbf p}_m\}_{m=1}^{\widetilde M}$，以降低场构建成本。对采样点 $\widetilde{\mathbf p}_m$，首先将其约束压力分配到最近的内部控制点：
\begin{equation}
    k^*(m)=\arg\min_{1\leq k\leq K}\|\widetilde{\mathbf p}_m-\mathbf{q}_k\|_2 .
    \label{eq:nearest-control}
\end{equation}
对第 $k$ 个控制点，其 CVF 方向由多个场分量叠加得到：
\begin{equation}
    \begin{aligned}
    \mathbf{F}_k =
    \sum_{c\in\mathcal{C}}
    \eta_c(s,V_c)\mathbf{F}_{k,c},\\
    \mathcal{C}=
    \{\mathrm{obs},\mathrm{nfz},\mathrm{risk},\mathrm{alt},\mathrm{curv},\mathrm{bnd}\},
    \end{aligned}
    \label{eq:cvf-field}
\end{equation}
其中 $\mathbf{F}_{k,c}$ 为第 $c$ 类约束或质量压力对应的方向分量，$\eta_c(s,V_c)$ 为与当前约束状态和局部违反程度有关的权重。
正式实验中，$(\mathrm{obs},\mathrm{nfz},\mathrm{risk},\mathrm{alt},\mathrm{curv},\mathrm{bnd})$ 的状态权重分别为：可行性形成 $(1.20,1.15,0.85,1.10,0.85,0.65)$，可行性保持 $(0.95,0.95,0.85,0.90,0.75,0.55)$，质量细化 $(0.35,0.35,0.55,0.40,0.95,0.35)$，停滞恢复 $(0.80,0.80,0.75,0.75,0.55,0.45)$。状态融合权重 $(\alpha_s,\beta_s,\gamma_s)$ 在保守调度下取 $(0.80,0.40,0.04)$、$(0.90,0.38,0.06)$、$(1.00,0.38,0.06)$ 和 $(0.90,0.28,0.04)$，分别对应上述四类状态。其中质量细化权重用于完整状态模板说明；在本文正式保守调度中，它主要作为候选权重定义保留，而非频繁触发的主导状态。

若多个采样点映射到同一内部控制点，先累加其压力再按命中次数求平均。记 $n_k$ 为映射到第 $k$ 个控制点且压力非零的采样点数量，则非曲率分量的控制点增量可写为
\begin{equation}
    \Delta \mathbf q_k^{\mathrm{path}}=
    \begin{cases}
    \dfrac{1}{n_k}\sum_{m:k^*(m)=k}\mathbf F(\widetilde{\mathbf p}_m), & n_k>0,\\[6pt]
    \mathbf 0, & n_k=0,
    \end{cases}
    \label{eq:cvf-average-pressure}
\end{equation}
其中 $\mathbf F(\widetilde{\mathbf p}_m)$ 为该采样点处障碍物、禁飞区、风险、高度和边界压力的加权和。曲率分量直接在控制点序列上计算，并与 $\Delta \mathbf q_k^{\mathrm{path}}$ 相加得到最终控制点增量 $\Delta\mathbf q_k$。

\paragraph{障碍物分量}
对于轴对齐障碍物 $\mathcal B=[x_1,x_2]\times[y_1,y_2]\times[z_1,z_2]$，令 $\mathbf c_{\mathcal B}(\mathbf p)$ 为 $\mathbf p$ 在该盒体上的最近点。若 $\mathbf p$ 位于障碍物外部，定义
\begin{equation}
    \mathbf d_o=
    [p_x-c_x,\ p_y-c_y,\ 0.35(p_z-c_z)]^\mathsf T,\qquad
    d_o=\|\mathbf d_o\|_2 .
    \label{eq:obstacle-distance-vector}
\end{equation}
当 $0<d_o\leq r_o$ 时，障碍物压力为
\begin{equation}
    \mathbf f_{\mathrm{obs}}(\mathbf p;\mathcal B)=
    \frac{r_o-d_o}{r_o}\frac{\mathbf d_o}{d_o}.
    \label{eq:obstacle-buffer-pressure}
\end{equation}
若 $\mathbf p$ 位于障碍物内部，则选择最近的水平外侧面法向 $\mathbf n_{\min}\in\{[-1,0,0]^\mathsf T,[1,0,0]^\mathsf T,[0,-1,0]^\mathsf T,[0,1,0]^\mathsf T\}$，并令
\begin{equation}
    \mathbf f_{\mathrm{obs}}(\mathbf p;\mathcal B)=
    (r_o+d_{\min})\mathbf n_{\min},
    \label{eq:obstacle-inside-pressure}
\end{equation}
其中 $d_{\min}$ 为到最近水平面的穿透深度，$r_o=7$ 为障碍物影响半径。垂向距离在式 \eqref{eq:obstacle-distance-vector} 中被压缩，是为了降低通过大幅爬升绕开水平约束的倾向。

\paragraph{禁飞区分量}
设柱状禁飞区为 $\mathcal Z=(\mathbf c_z,r_z,z_l,z_u)$。若 $p_z\notin [z_l,z_u]$，该分量为零；否则令 $\mathbf d_z=[p_x-c_{z,x},p_y-c_{z,y},0]^\mathsf T$，$d_z=\|\mathbf d_z\|_2$。禁飞区径向压力定义为
\begin{equation}
    \mathbf f_{\mathrm{nfz}}(\mathbf p;\mathcal Z)=
    \begin{cases}
    (r_z-d_z+r_n)\dfrac{\mathbf d_z}{d_z}, & d_z\leq r_z,\\[6pt]
    \dfrac{r_z+r_n-d_z}{r_n}\dfrac{\mathbf d_z}{d_z}, & r_z<d_z\leq r_z+r_n,\\[6pt]
    \mathbf 0, & d_z>r_z+r_n,
    \end{cases}
    \label{eq:nfz-pressure}
\end{equation}
其中 $r_n=7$ 为禁飞区缓冲影响半径；当 $d_z$ 接近零时使用固定水平单位方向避免除零。

\paragraph{风险分量}
风险分量不作为硬约束，而是质量压力。对风险热点 $\mathcal H_h=(\mathbf c_h,\sigma_h,a_h)$，令 $\widetilde{\mathbf d}_h=[p_x-c_{h,x},p_y-c_{h,y},0.35(p_z-c_{h,z})]^\mathsf T$，并定义
\begin{equation}
    \phi_h(\mathbf p)=
    a_h\exp\left(-\frac{\|\widetilde{\mathbf d}_h\|_2^2}{2\sigma_h^2}\right).
    \label{eq:risk-hotspot}
\end{equation}
仅当 $\phi_h(\mathbf p)>\phi_0$ 时触发风险压力：
\begin{equation}
    \mathbf f_{\mathrm{risk}}(\mathbf p;\mathcal H_h)=
    \kappa_r(\phi_h(\mathbf p)-\phi_0)
    \frac{\widetilde{\mathbf d}_h}{\|\widetilde{\mathbf d}_h\|_2+\epsilon_l},
    \label{eq:risk-pressure}
\end{equation}
其中 $\phi_0=0.24$，$\kappa_r=1.8$。该分量只降低风险暴露倾向，不改变障碍物、禁飞区和高度约束的可行性定义。

\paragraph{高度和边界分量}
记 $z_l=z_{\min}$，$z_u=z_{\max}$。高度分量采用分段形式：
\begin{equation}
f_{\mathrm{alt},z}(\mathbf p)=
\begin{cases}
z_l-p_z, & p_z<z_l,\\
z_u-p_z, & p_z>z_u,\\
0.35(z_l+b_z-p_z), & z_l\leq p_z<z_l+b_z,\\
-0.35(p_z-z_u+b_z), & z_u-b_z<p_z\leq z_u,\\
0, & \mathrm{otherwise},
\end{cases}
\label{eq:altitude-pressure}
\end{equation}
并叠加弱参考高度项 $0.04(z_{\mathrm{ref}}-p_z)$，其中 $b_z=3$。边界分量只在水平安全带内生效：
\begin{equation}
\begin{aligned}
f_{\mathrm{bnd},x}(\mathbf p)&=
\begin{cases}
x_{\min}+b_b-p_x, & p_x<x_{\min}+b_b,\\
x_{\max}-b_b-p_x, & p_x>x_{\max}-b_b,\\
0, & \mathrm{otherwise},
\end{cases}\\
f_{\mathrm{bnd},y}(\mathbf p)&\ \text{is defined analogously},
\end{aligned}
\label{eq:boundary-pressure}
\end{equation}
其中 $b_b=6$。该分量提高水平边界裕度，但边界硬约束仍由控制点投影保证。

\paragraph{转弯分量}
若连续控制点形成的局部转角 $\vartheta_k$ 接近或超过允许阈值，转弯分量对相关内部控制点施加局部平滑方向。$\vartheta_k$ 按式 \eqref{eq:turning-angle} 的夹角定义计算，但将采样点 $\mathbf p_m$ 替换为控制点 $\mathbf q_k$：
\begin{equation}
    \mathbf{F}_{k,\mathrm{curv}}=
    \begin{cases}
    \gamma_{\mathrm{curv}}
    \left[\tfrac{1}{2}(\mathbf{q}_{k-1}+\mathbf{q}_{k+1})-\mathbf{q}_k\right],
    & \vartheta_k\geq \theta_c,\\[4pt]
    \mathbf 0,
    & \mathrm{otherwise},
    \end{cases}
    \label{eq:curvature-field}
\end{equation}
其中 $\gamma_{\mathrm{curv}}=0.22$，$\theta_c=\alpha_{\mathrm{curv}}\theta_{\max}$，$\alpha_{\mathrm{curv}}=0.85$。该分量不会直接替代 B 样条平滑，而是降低控制多边形的局部尖锐程度；随后它还会经过式 \eqref{eq:cvf-field} 中的状态权重、式 \eqref{eq:component-clip} 中的逐维截断和式 \eqref{eq:norm-clip} 中的范数截断。

将全部内部控制点增量向量化后得到未截断 CVF 步长：
\begin{equation}
    \widetilde{\mathbf D}_{\mathrm{CVF}}=
    \mathrm{vec}
    \left(
    [\Delta\mathbf q_1^\mathsf{T},\ldots,\Delta\mathbf q_K^\mathsf{T}]^\mathsf{T}
    \right).
    \label{eq:cvf-vectorization}
\end{equation}
为避免 CVF 方向过大导致路径跨越可行走廊，所有分量叠加后需要进行逐维截断、强度缩放和整体范数截断：
\begin{equation}
    \bar D_{\mathrm{CVF},j}=
    \mathrm{clip}
    \left(
    \widetilde D_{\mathrm{CVF},j},
    -\rho_{\mathrm{step}}(u_j-l_j),
    \rho_{\mathrm{step}}(u_j-l_j)
    \right),
    \label{eq:component-clip}
\end{equation}
\begin{equation}
    \mathbf D_{\mathrm{CVF}}^{(0)}=
    \kappa_{\mathrm{cvf}}\bar{\mathbf D}_{\mathrm{CVF}},
    \label{eq:cvf-strength}
\end{equation}
\begin{equation}
    \|\mathbf{D}_{\mathrm{CVF}}\|_2
    \leq
    \rho_{\mathrm{norm}}\|\mathbf{u}-\mathbf{l}\|_2 ,
    \label{eq:norm-clip}
\end{equation}
其中 $\kappa_{\mathrm{cvf}}$ 为 CVF 强度系数。若 $\|\mathbf D_{\mathrm{CVF}}^{(0)}\|_2$ 超过式 \eqref{eq:norm-clip} 的上界，则按比例缩放到该上界内；否则令 $\mathbf D_{\mathrm{CVF}}=\mathbf D_{\mathrm{CVF}}^{(0)}$。其中 $\rho_{\mathrm{step}}$ 和 $\rho_{\mathrm{norm}}$ 为小比例截断系数。该有界设计是 CVF-AE 区别于普通势场排斥的关键：CVF 只提供受控的可行性方向，不允许单次候选生成主导整个搜索步长。

质量引导方向 $\mathbf D_{\mathrm Q}$ 由当前最优解、采样精英个体和控制点平滑项共同构成：
\begin{equation}
    \mathbf D_{\mathrm Q}=
    \mathrm{clip}
    \left[
    0.6(\mathbf X^*-\mathbf X_i)
    +0.4(\mathbf X_{\mathrm{elite}}-\mathbf X_i)
    +\mathbf D_{\mathrm{sm}},
    \rho_{\mathrm Q}
    \right],
    \label{eq:quality-step}
\end{equation}
其中 $\mathbf D_{\mathrm{sm}}$ 为控制点局部平滑增量，$\rho_{\mathrm Q}$ 表示逐维质量步长上界。

\subsection{稀疏触发与候选接收}

CVF-AE 仅在有必要时生成并评价 CVF 候选。触发策略综合考虑个体排名、是否接近可行、当前约束状态和代内评价预算。设第 $t$ 代状态为 $s_t$，对应触发比例为 $r_{s_t}$，则单代 CVF 候选数满足
\begin{equation}
    N_{\mathrm{cvf}}^{(t)}\leq
    \left\lceil r_{s_t}N\right\rceil .
    \label{eq:cvf-quota}
\end{equation}
对于可行性形成状态，算法优先处理违反量较小、具有进入可行区域潜力的个体；对于质量细化状态，算法减少 CVF 触发频率，重点保护已有可行精英；对于停滞恢复状态，算法允许少量额外 CVF 候选用于恢复搜索活性。

每个被触发的个体先评价基础候选 $\mathbf{X}_{\mathrm{AE}}$。若生成了 $\mathbf{X}_{\mathrm{CVF}}$，则对二者进行局部候选比较。CVF 候选被采用需同时满足两个条件：
\begin{equation}
    \mathbf{X}_{\mathrm{CVF}}
    \prec_{\mathrm{Deb}}
    \mathbf{X}_{\mathrm{AE}},
    \qquad
    J(\mathbf{X}_{\mathrm{CVF}})
    \leq
    J(\mathbf{X}_{\mathrm{AE}}),
    \label{eq:local-cvf-accept}
\end{equation}
其中 $\prec_{\mathrm{Deb}}$ 表示第 3 节定义的可行性优先比较关系。这是一个保守局部接收规则：CVF 候选不仅需要在可行性优先比较中优于基础候选，还不能带来更高综合目标值。该设计降低了 CVF 分支在早期以较高目标值强行替代 AE 候选的风险，但也意味着它不是最激进的可行性优先接收。若改为“可行性改善即接收”或允许不可行阶段目标容忍，应作为新的算法变体重新运行正式实验。最终候选再与父代个体按同一可行性优先规则比较，决定是否进入下一代。

算法还保留稀疏局部修复机制，用于对少量候选进行碰撞、禁飞区、转弯和风险方向上的局部校正。修复发生在候选生成之后，且不改变 CVF 的定位：CVF 是评价前的方向场引导，修复是额外的局部后处理，两者在功能上相互独立。

\subsection{算法复杂度}

设种群规模为 $N$，最大迭代次数为 $T$，路径采样点数量为 $M$，内部控制点数量为 $K$，空间约束元素数量为 $G$。基础进化搜索每代需要约 $N$ 次候选评价，总评价规模为 $O(TN)$。单次路径评价需要生成 B 样条采样点并检查目标项与约束项，其主要代价可概括为
\begin{equation}
    O\big(M(G+K)\big).
\end{equation}

CVF 会带来额外候选评价和场构造开销。若每代触发 CVF 的个体数为 $N_{\mathrm{cvf}}$，则额外评价规模为 $O(TN_{\mathrm{cvf}})$，场构造开销约为
\begin{equation}
    O\big(TN_{\mathrm{cvf}}M(G+K)\big).
\end{equation}
由于 $N_{\mathrm{cvf}}$ 受到稀疏触发比例和后期节流机制限制，通常满足 $N_{\mathrm{cvf}}\ll N$。因此，CVF-AE 的额外评价开销主要由少量关键个体承担，而不是对全体种群施加强制场更新。该设计使方法能够在强约束场景中提高可行路径形成能力，同时避免因全种群 CVF 更新带来评价预算膨胀。

\begin{algorithm}[htbp]
\caption{CVF-AE 的完整算法流程}
\label{alg:cvf-ae}
\begin{algorithmic}[1]
\Require 地图 $\mathcal{M}$，起点 $\mathbf{s}$，终点 $\mathbf{g}$，种群规模 $N$，最大迭代次数 $T$
\Ensure 最优可行路径 $\mathcal{P}^*$
\State 初始化种群 $\{\mathbf{X}_i\}_{i=1}^{N}$，并将控制点限制在边界内
\State 评价所有个体的 $J(\mathbf{X}_i)$ 与 $V(\mathbf{X}_i)$
\State 按 Deb 可行性优先规则记录当前最优解 $\mathbf{X}^*$
\For{$t=1$ to $T$}
    \State 根据可行率、平均违反量、多样性和近期改进识别约束状态 $s_t$
    \For{$i=1$ to $N$}
        \State 由基础自适应进化算子生成 $\mathbf{X}_{\mathrm{AE}}$
        \State 评价 $\mathbf{X}_{\mathrm{AE}}$ 的目标值和违反量
        \If{个体 $i$ 满足 CVF 稀疏触发条件}
            \State 构造 $\mathbf{D}_{\mathrm{CVF}}$ 与 $\mathbf{D}_{\mathrm{Q}}$
            \State 按式 \eqref{eq:cvf-candidate} 生成 $\mathbf{X}_{\mathrm{CVF}}$
            \State 评价 $\mathbf{X}_{\mathrm{CVF}}$ 的目标值和违反量
            \If{$\mathbf{X}_{\mathrm{CVF}} \prec_{\mathrm{Deb}} \mathbf{X}_{\mathrm{AE}}$ 且 $J(\mathbf{X}_{\mathrm{CVF}})\leq J(\mathbf{X}_{\mathrm{AE}})$}
                \State $\mathbf{Y}\gets \mathbf{X}_{\mathrm{CVF}}$
            \Else
                \State $\mathbf{Y}\gets \mathbf{X}_{\mathrm{AE}}$
            \EndIf
        \Else
            \State $\mathbf{Y}\gets \mathbf{X}_{\mathrm{AE}}$
        \EndIf
        \If{$\mathbf{Y}$ 满足稀疏修复条件}
            \State 对 $\mathbf{Y}$ 进行局部路径修复并重新评价
        \EndIf
        \If{$\mathbf{Y} \prec_{\mathrm{Deb}} \mathbf{X}_i$}
            \State $\mathbf{X}_i\gets \mathbf{Y}$
        \EndIf
        \State 更新当前最优解 $\mathbf{X}^*$
    \EndFor
\EndFor
\State \Return 由 $\mathbf{X}^*$ 解码得到的路径 $\mathcal{P}^*$
\end{algorithmic}
\end{algorithm}
